\section{Nombres premiers}

Les nombres premiers sont une notion de base en arithmétique. Ils forment en quelque sorte les \guillemets{briques} qui permettent de construire, par multiplication, l'ensemble de tous les nombres naturels. 




\medskip



\begin{definition}
On dit qu'un nombre positif $n$ est \textbf{premier} s'il possède exactement deux diviseurs.

\medskip
En d'autres mots, ce sont tous les nombres entiers $\geqslant 2$ et divisibles uniquement par 1 et lui-même.
\end{definition}

\begin{exemple}
2 ; 3 ; 5 ; 7 ; 11 sont des nombres premiers alors que 4 ; 6 ; 8 ; 9 ne le sont pas.
\end{exemple}


\smallskip
Nous savons depuis l'antiquité (et le génial Euclide) qu'il existe une infinité de nombres premiers. Le \textit{théorème fondamental de l'arithmétique} assure de plus que tout nombre entier se factorise de manière unique (à l'ordre près) en un produit de nombres premiers.

\smallskip
Il est intéressant de remarquer qu'il n'existe à ce jour aucune manière \guillemets{efficace} de savoir si un nombre est premier\footnote{On parle bien entendu ici de nombres comportant de milliers, voir millions de chiffres.}. C'est cette faille dans nos les connaissances arithmétiques actuelles qui a permis l'élaboration de systèmes cryptographiques utilisés dans le monde entier pour coder une information numérique. 

\smallskip
Bien que de prime abord les nombres entiers semblent être des objets simples, bon nombre de problèmes les concernant n'ont pas encore été résolus.


\begin{lemme}\label{lem:divpremier}
Tout entier $n \geqslant 2$ admet un diviseur qui est un nombre premier.
\end{lemme}

\noindent{\it Démonstration:}

Soit $\mathcal{D}$ l'ensemble des diviseurs de $n$ qui sont $\geqslant 2$ :
$$\mathcal{D} = \big\{ k \geqslant 2 \mid \  k | n \big\}.$$
L'ensemble $\mathcal{D}$ est non vide (car $n \in \mathcal{D}$) et $\mathcal{D}$ est fini (car il ne peut pas contenir plus de $n$ éléments).

Donc $\mathcal{D}$ possède un plus petit élément qu'on note $p = \min \mathcal{D}$.

\smallskip
Supposons, par l'absurde, que $p$ ne soit pas un nombre premier. Alors $p$ admet un diviseur
$q$ tel que $1 < q < p$. Mais si $q|p$ et $p|n$ alors $q|n$. Donc $q\in \mathcal{D}$ et $q<p$ ce qui contredit la minimalité de $p$.

\smallskip
Conclusion : $p$ est un nombre premier. Et comme $p \in \mathcal{D}$, $p$ divise $n$.

\hfill $\blacksquare$

\vspace{5mm}

\begin{proposition}
Il existe une infinité de nombres premiers.
\end{proposition}

\noindent{\it Démonstration:}

Par l'absurde, supposons qu'il n'y ait qu'un nombre fini de nombres premiers que l'on note
$p_1=2$, $p_2=3$, $p_3$,\ldots, $p_n$ et considérons l'entier $N=p_1\cdot p_2\cdot \ldots \cdot p_n+ 1$.

\smallskip
Soit $p$ un diviseur premier de $N$ (un tel $p$ existe par le lemme précédent). D'une part, comme  $p$ est l'un des entiers $p_i$, $p | p_1\cdot \ldots \cdot p_n$ et d'autre part $p|N$. Donc $p$ divise la différence $N-p_1\cdot \ldots \cdot p_n=1$.
Cela implique que $p=1$, ce qui contredit que $p$ soit un nombre premier.

Cette contradiction nous permet de conclure qu'il existe une infinité de nombres premiers.

\hfill $\blacksquare$


\newcommand{\entourer}[1]{\ovalbox{\color{blue}#1}}
\newcommand{\barrer}[1]{{\cancel{\color{gray} #1}}}
\centerline{\textsc{Le crible d'Eratosthène}}

\bigskip
Le \textbs{crible d'Eratosthène} est un algorithme qui permet de trouver les premiers nombres premiers.
On commence par écrire une liste avec les premiers entiers : pour notre exemple de $2$ à $25$.
$$2\ \ \ 3\ \ \ 4\ \ \ 5\ \ \ 6\ \ 7\ \ \ 8\ \ \ 9\ \ \ 10\ \ \ 11\ \ \ 12\ \ \ 13\ \ \ 14\ \ 15\ \
\ 16\ \ \ 17\ \ \ 18\ \ \ 19\ \ \ 20\ \ \ 21\ \ \ 22\ \ \ 23\ \ \ 24\ \ \ 25$$
Nous savons que $2$ est le premier nombre premier.
On entoure $2$. Ensuite on barre
 tous les multiples suivants de $2$ qui ne seront donc pas premiers
(car divisible par $2$) :
$$\entourer{2}\ \ 3\ \ \barrer{4}\ \ 5\ \ \barrer{6}\ \ 7\ \ \barrer{8}\ \ 9\ \ \barrer{10}\ \ 11\ \ \barrer{12}\ \
13\ \ \barrer{14}\ \ 15\ \
\barrer{16}\ \ 17\ \ \barrer{18}\ \ 19\ \ \barrer{20}\ \ 21\ \ \barrer{22}\ \ 23\ \ \barrer{24}\ \ 25$$

Le premier nombre suivant et non-barré de la liste est $3$. Celui-ci est nécessairement premier car il n'est pas un multiple des nombres qui lui sont inférieurs et donc pas divisible par eux non plus) (sinon il serait rayé). On entoure $3$ et on raye tous les multiples suivants de $3$ :
$$\entourer{2}\ \ \entourer{3}\ \ \barrer{4}\ \ 5\ \ \barrer{6}\ \ 7\ \ \barrer{8}\ \ \barrer{9}\ \ \barrer{10}\ \ 11\ \ \barrer{12}\ \
13\ \ \barrer{14}\ \ \barrer{15}\ \
\barrer{16}\ \ 17\ \ \barrer{18}\ \ 19\ \ \barrer{20}\ \ \barrer{21}\ \ \barrer{22}\ \ 23\ \ \barrer{24}\ \ 25$$
%Le premier nombre restant est $5$ et est donc premier. On raye les multiples de $5$.
%$$\entourer{2}\ \ \entourer{3}\ \ \barrer{4}\ \ \entourer{5}\ \ \barrer{6}\ \ 7\ \ \barrer{8}\ \ \barrer{9}\ \ \barrer{10}\ \ 11\ \ \barrer{12}\ \
%13\ \ \barrer{14}\ \ \barrer{15}\ \
%\barrer{16}\ \ 17\ \ \barrer{18}\ \ 19\ \ \barrer{20}\ \ \barrer{21}\ \ \barrer{22}\ \ 23\ \ \barrer{24}\ \ \barrer{25}$$
%$7$ est donc premier, on raye les multiples de $7$ (ici pas de nouveaux nombres à barrer).
%Ainsi de suite : $11, 13, 17, 19, 23$ sont premiers.
On continue le processus ainsi de suite. Les nombres encadrés sont les premiers nombres premiers :
$$\entourer{2}\ \ \entourer{3}\ \ \barrer{4}\ \ \entourer{5}\ \ \barrer{6}\ \ \entourer{7}\ \ \barrer{8}\ \
\barrer{9}\ \ \barrer{10}\ \ \entourer{11}\ \ \barrer{12}\ \
\entourer{13}\ \ \barrer{14}\ \ \barrer{15}\ \
\barrer{16}\ \ \entourer{17}\ \ \barrer{18}\ \ \entourer{19}\ \ \barrer{20}\ \ \barrer{21}\ \ \barrer{22}\ \ \entourer{23}
%\ \ \barrer{24}\ \ \barrer{25}
$$



\begin{theoreme}[\normalsize\textit{théorème fondamental de l'arithmétique}]\label{thm:TFA}
Soit $n \ge 2$ un entier. Il existe des nombres premiers $p_1 < p_2 < \cdots < p_r$
et des exposants entiers $\alpha_1, \alpha_2, \dots, \alpha_r \ge 1$ tels que :
$$n = p_1^{\alpha_1} \times p_2^{\alpha_2} \times \cdots \times p_r^{\alpha_r}.$$
De plus les $p_i$ et les $\alpha_i$ ($i=1,\ldots,r$) sont uniques.
\end{theoreme}


\vspace{5mm}
\begin{remarque}
Voici la raison pour laquelle on ne veut pas que $1$ soit un nombre premier. Si c'était le cas, il y aurait une infinité de façons de décomposer un nombre : $$24 = 2^3 \times 3 = 1 \times 2^3 \times 3 =
1^2 \times 2^3 \times 3 = \cdots$$
\end{remarque}



\noindent{\it Démonstration:}

\medskip
\textbf{Existence.}

\smallskip
Nous allons démontrer l'existence de la décomposition par récurrence sur l'entier $n$.

L'entier $n=2$ est déjà décomposé. Soit $n \geqslant 3$ et supposons que
tout entier $< n$ admette une décomposition en facteurs premiers.

\smallskip
Notons $p_1$ le plus petit nombre premier divisant $n$ (voir le lemme \ref{lem:divpremier}).
Si $n$ est un nombre premier alors $n=p_1$ et c'est fini. Sinon
on définit l'entier $n'= \frac{n}{p_1} < n$ et on applique notre hypothèse de récurrence
à $n'$ qui admet une décomposition en facteurs premiers. Alors $n = p_ 1 \times n'$ admet aussi
une décomposition.

\bigskip

%\textbf{Unicité.}
%
%\smallskip
%Nous allons démontrer qu'une telle décomposition est unique en effectuant cette fois
%une récurrence sur la somme des exposants $\sigma = \sum_{i=1}^{r} \alpha_i$.
%
%Si $\sigma=1$ cela signifie $n=p_1$ qui est bien l'unique écriture possible.
%
%Soit $\sigma \ge 2$. On suppose que les entiers dont la somme des exposants est $< \sigma$
%ont une unique décomposition. Soit $n$ un entier dont la somme des exposants vaut
%$\sigma$. \'Ecrivons le avec deux décompositions :
%$$n= p_1^{\alpha_1} \times p_2^{\alpha_2} \times \cdots \times p_r^{\alpha_r}
% = q_1^{\beta_1} \times q_2^{\beta_2} \times \cdots \times q_s^{\beta_s}.$$
%Sans perte de généralité, on peut supposer qu'on a $p_1< p_2 < \ldots$ et $q_1 < q_2 < \ldots$.
%
%\bigskip
%Si $p_1 < q_1$ alors $p_1 < q_j$ pour tous les $j=1,\ldots,s$. Ainsi
%$p_1$ divise $p_1^{\alpha_1} \times p_2^{\alpha_2} \times \cdots \times p_r^{\alpha_r} = n$
%mais ne divise pas $q_1^{\beta_1} \times q_2^{\beta_2} \times \cdots \times q_s^{\beta_s} = n$. Ce qui est absurde.
%Donc $p_1 \ge q_1$.
%
%\bigskip
%Si $p_1 > q_1$ le même raisonnement conduit aussi à une contradiction. On conclut que $p_1=q_1$.
%
%\bigskip
%On pose alors
%$$n' = \frac{n}{p_1} = \frac{n}{q_1} = p_1^{\alpha_1-1} \times p_2^{\alpha_2} \times \cdots \times p_r^{\alpha_r}
% = q_1^{\beta_1-1} \times q_2^{\beta_2} \times \cdots \times q_s^{\beta_s}$$
%L'hypothèse de récurrence qui s'applique à $n'$ implique que ces deux décompositions sont les mêmes.
%Ainsi $r=s$ et $p_i=q_i$, $\alpha_i = \beta_i$ pour tout $i=1,\ldots,r$.



\bigskip

\begin{proposition}[Lemme d'Euclide]
\index{lemme!d'Euclide}
Soit $p$ un nombre premier.
Si $p | ab$ alors $p|a$ ou $p | b$.
\end{proposition}

\begin{demonstration}
Si $p$ ne divise pas $a$ alors $p$ et $a$ sont premiers entre eux
(en effet les diviseurs de $p$ sont $1$ et $p$, mais seul $1$ divise aussi $a$, donc $\pgdc(a,p)=1$).
Ainsi par le lemme de Gauss $p | b$.
\end{demonstration}

\bigskip
%EXERCICES (voir théorie arithmétique 19-20)


\begin{exstheorie}
\exo
Démontrer que si un nombre entier $n$ peut s'écrire comme le produit de 2 nombres entiers $a$ et $b$ ($n=a\cdot b$) alors $a\leqslant \sqrt{n}$ ou $b\leqslant\sqrt{n}$.

\bigskip
\textit{\small\textbf{Remarque :}
Il découle de cet exercice que, pour déterminer si un nombre est premier par divisions successives des entiers qui lui sont inférieurs, il suffit de tester jusqu'à $\sqrt{n}$. De plus il suffit de tester uniquement les nombres premiers.\\
En particulier, pour déterminer si un nombre $n\leqslant 100$ est premier, il suffit de tester la divisibilité de $n$ par $2$, $3$, $5$ et $7$.
}




\eexo





\exo
Démonter que si $p$ est un nombre premier, alors $\sqrt{p}$ est irrationnel.

\textit{\small\textbf{Remarque :}
Raisonner par l'absurde et utiliser le lemme d'Euclide.}




\eexo





\exo
On souhaite créer une fonction en Python qui simule le \textit{Crible d'Eratosthène}. On utilisera un $0$ pour représenter les nombres censés être barrés dans la liste puis on veillera à supprimer ces zéros pour que la fonction renvoie au final une liste ne contenant que les nombres premiers.

\medskip
Dans un en premier temps, il s'agit de générer une liste ``\comp{liste\_entiers}'' qui contient tous les entiers de $0$ à $n$ à l'aide de l'instruction Python

\begin{center}
\comp{liste\_entiers = list(range(n+1))}
\end{center}

Remarquons que chaque entier de la liste est égal à sa position (indice) dans la liste. Les étapes à réaliser sont les suivantes :

\begin{enumerate}[label=\roman*),itemsep=0mm,left=-2mm]
\item Une boucle d'indice $k$ qui va parcourir les positions (nombres entiers) de la liste.
\item Vérifier si la position contient un ``nombre non-barré''.
\item A l'aide d'une deuxième boucle (imbriquée dans la première), parcourir tous les indices qui sont des multiples du nombre $k$ (et plus grands que $k$) et leur attribuer la valeur $0$.
\item Utiliser les commandes spécifiques aux listes pour enlever les $0$ de la liste.
\end{enumerate}

\end{exstheorie}


















